% Research candidate, not yet inserted into the main manuscript.
% Relies on Lemma~\ref{lem:reduction} and Proposition~\ref{prop:potential}.
\subsection*{A first-passage and residual decomposition (research candidate)}
\label{subsec:residual}
Fix an odd $c\ge-1$ with $3\nmid c$ and set
$\gamma_0=2/\log(4/3)$. Fix $0<\beta<1/3$ and $X>1$.
Let $W_X=(X^{1-\beta},X]\cap\mathbb N$ and
$X_1=X^{1-2\beta}$. For an odd starting value $m$, define
$\Phi_X(m)$ as in the proof of Theorem~\ref{thm:lower}
with the residue restriction removed: it is the sum of those values
in $W_X$ occurring in the Syracuse blocks up to and including the
first Syracuse passage below $X_1$, and is defined to be zero if
that first passage never occurs.

\begin{proposition}[Exact first-passage and residual decomposition]
\label{prop:residual_decomposition}
Assume the predecessor-potential statement of
Proposition~\ref{prop:potential}(a) and the window-reduction
Lemma~\ref{lem:reduction}. For each fixed $X$, the logarithmic mean
\[
 F_X=\lim_{N\to\infty}\frac1{\log N}
 \sum_{1\le w<N}\frac{\Phi_X(w^\flat)}{w}
\]
exists. Moreover the nonnegative residual mean
\[
 R_X=\lim_{N\to\infty}\frac1{\log N}
 \sum_{1\le w<N}\frac1w
 \left(\sum_{\substack{a\le X\\a\in\operatorname{orb}(w)}}a
       -\Phi_X(w^\flat)\right)
\]
exists and satisfies the exact identity
\[
 \boxed{\sum_{a\le X}a\,h(a)=F_X+R_X,\qquad R_X\ge0.}
\]
In particular, if $F_X=\gamma_0X+o(X)$, then the remaining
upper-bound conjecture
\[
 \limsup_{X\to\infty}\frac1X\sum_{a\le X}a\,h(a)
 \le\gamma_0
\]
is equivalent to $R_X=o(X)$. Without this additional calibration,
the two sufficient tasks are
$\limsup F_X/X\le\gamma_0$ and $R_X=o(X)$.
\end{proposition}

\begin{proof}
By the distinctness argument in Step~1 of the proof of
Theorem~\ref{thm:lower}, the blocks counted by $\Phi_X(m)$
contain no repeated positive integer. Thus
$0\le\Phi_X(m)\le B_X:=\sum_{a\in W_X}a$.
Choose $x_0>X$ large enough that
$\tfrac34 x\exp(-(\log x)^{0.7})>X$ for every $x\ge x_0$.
If $m>x$ is odd, its first passage $M=\operatorname{Pass}_x(m)$
exists, and $M>X$, then the initial Syracuse blocks before $M$
contain no value at most $X$, and the later blocks are precisely the
Syracuse blocks of $M$. Hence $\Phi_X(m)=\Phi_X(M)$.
Consequently $\Phi_X/B_X$ satisfies
Lemma~\ref{lem:reduction}, which yields the existence of $F_X$.
By Proposition~\ref{prop:potential}(a), the logarithmic mean of
$\sum_{a\le X,\,a\in\operatorname{orb}(w)}a$ exists and equals
$\sum_{a\le X}a\,h(a)$, by finite linearity. The pointwise
inequality from Step~1 of Theorem~\ref{thm:lower}
shows that the difference inside the definition of $R_X$ is
nonnegative. Subtraction of the two convergent logarithmic means
proves the identity. The final equivalence follows from
$R_X\ge0$ and the assumed asymptotic for $F_X$.
\end{proof}

\paragraph{Research implications.}
The proof of Theorem~\ref{thm:lower} already supplies
$F_X\ge\gamma_0X-o(X)$, by its truncated first-passage comparison.
Two genuinely additional estimates are required for the matching
upper bound: a reverse (upper) first-passage calibration for $F_X$,
and a bound $R_X=o(X)$ on the residual orbit mass.
The existing coupling proves a lower bound without having to
control $\Phi_X$ on the exceptional coupling event. An upper
bound must additionally address that event, for example by a
uniform-integrability argument. The factor $B_X=O(X^2)$ prevents
small exceptional-event probability alone from being sufficient.
